\chapter{Pipeline Stage Reference}
\label{app:stages}

Chapter~\ref{ch:design} describes what each stage does and why the architecture
takes the shape it does. This appendix records the file-level interface: what
each stage reads and what it writes; the mapping from an analysis to its
inference status is in the analysis register (Appendix~\ref{sec:register})
rather than repeated here. Paths are relative to \texttt{sepsis/project/}; all intermediate tables
are Parquet, with a CSV mirror written for any result table small enough to read
directly.

\begin{table}[H]
    \centering\scriptsize\setlength{\tabcolsep}{4pt}
    \begin{tabularx}{\textwidth}{p{0.9cm}p{4.2cm}X}
        \toprule
        \textbf{Stage} & \textbf{Reads} & \textbf{Writes} \\
        \midrule
        01 & MIMIC-IV source tables &
        n/a (validation only; halts on a missing table) \\
        \addlinespace
        02 & \texttt{hosp/}, \texttt{icu/} source tables &
        \texttt{02\_cohort\_\{A..F\}}, \texttt{02\_label\_variance\_table},
        \texttt{02\_alt\_label\_variance\_table} \\
        \addlinespace
        03 & \texttt{02\_cohort\_*}; \texttt{chartevents}, \texttt{labevents},
        \texttt{inputevents} (streamed) &
        \texttt{03\_person\_hours\_\{A..F\}}, \texttt{03\_early\_onset\_shift} \\
        \addlinespace
        04 & \texttt{03\_person\_hours\_*} (three columns) &
        \texttt{04\_sample\_size} \\
        \addlinespace
        05 & \texttt{03\_person\_hours\_*} &
        \texttt{05\_model\_primary\_*}, \texttt{05\_coef\_all\_variants},
        \texttt{05\_predictions\_primary\_*}, \texttt{05\_h2\_stability\_*};
        \texttt{*\_ablated} counterparts \\
        \addlinespace
        06 & \texttt{03\_person\_hours\_*} &
        \texttt{06\_model\_gbt\_*}, \texttt{06\_predictions\_comparators\_*},
        \texttt{06\_gbt\_importance\_*}, \texttt{06\_schoenfeld\_*};
        ablated GBT counterparts \\
        \addlinespace
        07 & \texttt{05\_*}, \texttt{06\_*} predictions &
        \texttt{07\_metric\_results}, \texttt{\_random}, \texttt{\_altlabel},
        \texttt{07\_utility\_sweep}, \texttt{07\_utility\_max},
        \texttt{07\_ablation\_metrics} \\
        \addlinespace
        08 & eICU-CRD source tables; frozen \texttt{05\_}/\texttt{06\_} models
        (including the Variant D fits, for the Amendment 9 arm) &
        \texttt{08\_external\_validation\_results}, \texttt{\_abxonly},
        \texttt{\_altlabel}, \texttt{08\_eicu\_predictions\_*} (A--C on the
        microbiology-covered sub-cohort, D on the full cohort),
        \texttt{08\_eicu\_coverage},
        \texttt{08\_event\_rate\_reconciliation}, \texttt{\_altlabel},
        \texttt{08\_h5\_internal\_external} \\
        \addlinespace
        09 & \texttt{07\_metric\_results*} &
        \texttt{09\_variance\_decomposition\_table},
        \texttt{09\_hypothesis\_verdicts},
        \texttt{09\_label\_anchor\_attribution},
        \texttt{09\_label\_agreement}, \texttt{09\_pretreatment\_onsets} \\
        \addlinespace
        10 & \texttt{03\_person\_hours\_*}, \texttt{05\_}/\texttt{06\_}
        predictions &
        Per-stratum performance and label-sensitivity tables \\
        \addlinespace
        11 & \texttt{05\_}/\texttt{06\_}/\texttt{08\_} predictions &
        \texttt{11\_calibration\_summary}, \texttt{11\_operating\_points},
        \texttt{11\_net\_benefit\_at\_prevalence},
        \texttt{11\_patient\_level\_metrics},
        \texttt{11\_lead\_time\_summary}, external counterparts \\
        \addlinespace
        12 & Stored predictions only (refits nothing) &
        \texttt{12\_confirmatory\_tests}, \texttt{12\_internal\_ci},
        \texttt{12\_subgroup\_auroc\_ci}, \texttt{\_contrasts},
        \texttt{12\_subgroup\_labelsens\_*}, \texttt{12\_altlabel\_*}
        (including \texttt{\_transport\_ci}),
        \texttt{12\_ablation\_ci}, \texttt{12\_calibration\_ci},
        \texttt{\_contrasts},
        \texttt{12\_h2\_*}, \texttt{12\_analysis\_register} \\
        \bottomrule
    \end{tabularx}
    \caption{Stage inputs and outputs. Which result each stage supports is
    recorded once, in the analysis register of
    Appendix~\ref{sec:register}, and is not repeated here. Stage 03 is
    the only stage that reads the full-size source tables, and Stage 12 is the
    only one that reads nothing but stored predictions, which is what allows
    the inference layer to be re-run and checked without refitting any model.}
    \label{tab:stage_reference}
\end{table}

\paragraph{Naming conventions.} A file's numeric prefix is the stage that wrote
it. A \texttt{\_\{A..F\}} suffix denotes the label variant; \texttt{\_ablated}
denotes the Protocol Amendment 5 feature arm; \texttt{\_random} denotes the
random-split arm used only for H4. Files without a variant suffix are pooled
across variants by construction. No stage overwrites a file written by another
stage.
